% LaTeX template for "A Traveler's Companion to the Mazes of Menace"
% 5" × 8" handbook format, used with pandoc and latex-filter.lua
\documentclass[openany,10pt]{book}

% --- Page geometry: A5 (148 × 210 mm = 5.83 × 8.27 in) with binding margins ---
\usepackage[
  paperwidth=148mm,
  paperheight=210mm,
  inner=0.6in,
  outer=0.5in,
  top=0.65in,
  bottom=0.6in,
  headheight=14pt,
  headsep=0.15in,
  footskip=0.3in
]{geometry}

% --- Symbol fonts and unicode glyph fallbacks ---
% EB Garamond has no ★ (U+2605) or ⅕ (U+2155); remap them to LaTeX
% equivalents so xelatex doesn't print .notdef boxes.
\usepackage{amsmath}
\usepackage{amssymb}
\usepackage{newunicodechar}
\newunicodechar{★}{\ensuremath{\bigstar}}
\newunicodechar{⅕}{\ensuremath{\tfrac{1}{5}}}
% APL QUAD (U+2395) — used as an empty-cursor box for the ghost glyph in the
% Special Symbols table. Source Code Pro doesn't ship it; render it manually
% as a thin framed box at line height.
\newunicodechar{⎕}{{\setlength{\fboxsep}{0pt}\framebox[0.55em]{\rule{0pt}{1ex}}}}

% --- Fonts via fontspec (XeLaTeX) ---
\usepackage{fontspec}
\setmainfont{EB Garamond}[
  Path=fonts/,
  Extension=.ttf,
  UprightFont=EBGaramond,
  ItalicFont=EBGaramond-Italic,
  BoldFont=EBGaramond-Bold,
  BoldItalicFont=EBGaramond-BoldItalic
]
\setmonofont{Source Code Pro}[
  Path=fonts/,
  Extension=.ttf,
  UprightFont=SourceCodePro-Medium,
  ItalicFont=SourceCodePro-MediumItalic,
  BoldFont=SourceCodePro-Bold,
  BoldItalicFont=SourceCodePro-BoldItalic,
  Scale=0.80
]

% --- Typography ---
\usepackage{microtype}
\usepackage{setspace}
\setstretch{1.1}

% --- No-indent paragraphs with inter-paragraph spacing ---
\setlength{\parindent}{0pt}
\setlength{\parskip}{0.4em plus 0.1em minus 0.05em}

% --- Colors ---
\usepackage{xcolor}
\definecolor{codeframe}{gray}{0.35}

% --- Etoolbox for conditional commands ---
\usepackage{etoolbox}

% --- Running part name for headers ---
\def\currentpartname{}

% --- Headers and footers ---
\usepackage{fancyhdr}
\pagestyle{fancy}
\fancyhf{}
\renewcommand{\headrulewidth}{0.3pt}
\renewcommand{\footrulewidth}{0pt}

% Redefine marks: chapter sets \leftmark, no section marks
\renewcommand{\chaptermark}[1]{\markboth{#1}{}}
\renewcommand{\sectionmark}[1]{}

% Running header: "Part Name · Chapter Name" on left, page on right
% Uses centered dot with generous space; omits dot when no chapter set
\fancyhead[L]{\small\itshape
  \ifdefempty{\currentpartname}%
    {\leftmark}%
    {\currentpartname\ifx\leftmark\empty\else\hspace{0.8em}\textperiodcentered\hspace{0.8em}\leftmark\fi}}
\fancyhead[R]{\small\thepage}

% Plain style for chapter openers: just page number, no rule
\fancypagestyle{plain}{
  \fancyhf{}
  \fancyhead[R]{\small\thepage}
  \renewcommand{\headrulewidth}{0pt}
}

% Part page style: completely empty
\fancypagestyle{partpage}{
  \fancyhf{}
  \renewcommand{\headrulewidth}{0pt}
}

% --- Chapter/section formatting ---
\usepackage{titlesec}
\usepackage{needspace}

% Chapter formatting. We override \@makechapterhead to drop book.cls's
% default 50pt + 40pt vertical padding (which pushed chapter content
% well below the page header), letting more material (e.g. the long
% wand table + first subsection) land on the chapter's first page.
\makeatletter
\def\@makechapterhead#1{%
  \vspace*{0pt}%
  {\parindent \z@ \raggedright \normalfont
   \LARGE\bfseries #1\par\nobreak
   \vskip 10pt}}
\makeatother
\titleformat{\chapter}[display]
  {\normalfont\LARGE\bfseries}
  {}
  {0pt}
  {\LARGE}
\titlespacing*{\chapter}{0pt}{0pt}{0.15in}

% Section formatting — needspace ensures heading stays with following content
\titleformat{\section}
  {\needspace{5\baselineskip}\normalfont\Large\bfseries}
  {}
  {0pt}
  {}
\titlespacing*{\section}{0pt}{1.5em}{0.5em}

% Subsection formatting. needspace bumped to 7 lines so a #### heading
% never strands at the bottom of a page above a table that has been
% pushed onto the next page (the Dart, Boots, and Liches headings
% were the recurring offenders).
\titleformat{\subsection}
  {\needspace{7\baselineskip}\normalfont\large\bfseries}
  {}
  {0pt}
  {}
\titlespacing*{\subsection}{0pt}{1.2em}{0.3em}

% Subsubsection (italic bold)
\titleformat{\subsubsection}
  {\needspace{3\baselineskip}\normalfont\normalsize\bfseries\itshape}
  {}
  {0pt}
  {}

% --- Remove chapter/section numbering ---
\setcounter{secnumdepth}{-2}

% --- Lists ---
\usepackage{enumitem}
\setlist{nosep, leftmargin=1.5em}
\setlist[itemize]{label={\raisebox{0.4ex}{\fontsize{3pt}{3pt}\selectfont◆}}}

% --- Links ---
\usepackage{hyperref}
\hypersetup{
  colorlinks=true,
  linkcolor=black,
  urlcolor=black,
  pdfborder={0 0 0},
  % Tell PDF viewers (Adobe Reader, Preview, etc.) that in two-page
  % spread mode the first page should sit on the right, leaving the
  % left blank. This makes facing-page view match the printed book,
  % where odd-numbered pages are rectos.
  pdfpagelayout=TwoPageRight
}

% --- Tables ---
\usepackage{longtable}
\usepackage{booktabs}
\usepackage{array}
\usepackage{calc}
\usepackage{multicol}  % two-column index appendix
\renewcommand{\arraystretch}{1.15}
% Pandoc uses \def\LTcaptype{none} for uncaptioned tables
\newcounter{none}
% Reduce table font size to fit 5-inch pages
\AtBeginEnvironment{longtable}{\footnotesize}

% Suppress phantom longtable headers. longtable's pagination
% sometimes emits a fresh \endhead at the top of the page after the
% table actually ended, leaving a ghost column header floating
% before the next section. Loading ltablex with \keepXColumns
% patches the longtable environment, the \null acts as a marker
% that the table is genuinely complete, and the trailing
% \vspace*{-\baselineskip} pulls back the visible empty row that
% the \null would otherwise leave.
%
% Separately, \LTpre and \LTpost (the glue before/after each
% longtable) default to \bigskipamount = 12pt, which combined with
% the surrounding paragraph skip leaves a chunk of empty space
% above and below every table. Tighten them so tables sit closer
% to the surrounding prose.
\usepackage{ltablex}
\keepXColumns
\setlength\LTpre{\smallskipamount}
\setlength\LTpost{\smallskipamount}
% The ghost-header marker goes AFTER \end{longtable} (via
% \AfterEndEnvironment) so it doesn't get treated as a phantom
% trailing row inside the table itself.
\AfterEndEnvironment{longtable}{\null\vspace*{-\baselineskip}}

% --- Code blocks ---
\usepackage{fancyvrb}
\usepackage{fvextra}

% Shaded environment: minipage to prevent page breaks in code
\newenvironment{Shaded}{%
  \begin{minipage}{\linewidth}\vspace{0.5em}%
}{%
  \vspace{0.5em}\end{minipage}%
}

% Pandoc highlighting macros (token styles, may redefine Shaded)
$if(highlighting-macros)$
$highlighting-macros$
% Re-override Shaded with our minipage version
\renewenvironment{Shaded}{%
  \begin{minipage}{\linewidth}\vspace{0.5em}%
}{%
  \vspace{0.5em}\end{minipage}%
}
$endif$

% Override Highlighting for our formatting
\DefineVerbatimEnvironment{Highlighting}{Verbatim}{
  fontsize=\footnotesize,
  baselinestretch=0.85,
  frame=single,
  framesep=0.3em,
  rulecolor=\color{codeframe},
  commandchars=\\\{\},
  breaklines=true,
  breaksymbol={}
}

% Plain verbatim: also footnotesize with frame.
% Tight baselinestretch so line-drawing chars (|, -) touch
% vertically in ASCII figures.
\DefineVerbatimEnvironment{verbatim}{Verbatim}{
  fontsize=\footnotesize,
  baselinestretch=0.85,
  breaklines=true,
  breakanywhere=true,
  frame=single,
  framesep=0.3em,
  rulecolor=\color{codeframe}
}

% Prevent plain verbatim from splitting across pages
\BeforeBeginEnvironment{verbatim}{\begin{minipage}{\linewidth}\vspace{0.3em}}
\AfterEndEnvironment{verbatim}{\vspace{0.3em}\end{minipage}}

% --- Block quotes ---
\usepackage{quoting}
\renewenvironment{quote}
  {\begin{quoting}[leftmargin=1.5em, rightmargin=1em, vskip=0.5em]\small\itshape}
  {\end{quoting}}

% --- Graphics ---
\usepackage{graphicx}
\usepackage{grffile}

% Pandoc 3.5+ wraps every \includegraphics in \pandocbounded, which
% scales the image down to fit within \linewidth x \textheight. We
% accept that behaviour but define it as a passthrough so older
% pandocs (or our own definition) keep working.
\providecommand{\pandocbounded}[1]{%
  \begingroup
    \setbox0=\hbox{#1}%
    \ifdim\wd0>\linewidth
      \resizebox{\linewidth}{!}{\box0}%
    \else
      \box0%
    \fi
  \endgroup
}

% --- Widows and orphans ---
\widowpenalty=10000
\clubpenalty=10000

% --- Tight list support for pandoc ---
\providecommand{\tightlist}{%
  \setlength{\itemsep}{0pt}\setlength{\parskip}{0pt}}

% --- Pandoc column width support ---
\makeatletter
\newlength{\currentparskip}
\providecommand{\columnbreak}{\vfill\break}
$if(tables)$
\usepackage{longtable,booktabs,array}
$if(multiline-tables)$
\usepackage{multirow}
$endif$
\providecommand{\toprule}[1][]{\hline}
\providecommand{\bottomrule}[1][]{\hline}
\providecommand{\midrule}[1][]{\hline}
$endif$
\makeatother

% --- CSL refs support (pandoc may generate these) ---
$if(csl-refs)$
\newlength{\cslhangindent}
\setlength{\cslhangindent}{1.5em}
\newenvironment{CSLReferences}[2]{}{}
$endif$

% ============================================================
\begin{document}

% --- Quick Command Reference (inside-cover cheat sheet) ---
\thispagestyle{empty}
{\centering{\large\bfseries Quick Command Reference}\par\vspace{0.2em}
{\itshape The keys you reach for most. There are many more; press \texttt{?} in game for the full list.}\par}
\vspace{1.2em}
\footnotesize
\noindent
\begin{minipage}[t]{0.48\textwidth}
\textbf{Getting around}\par\vspace{0.2em}
\begin{tabular}{@{}ll@{}}
\texttt{h j k l} & move west, south, north, east\\
\texttt{y u b n} & move diagonally\\
\texttt{H J K L\ldots} & run until something happens\\
\texttt{\_} & travel to a chosen spot\\
\texttt{<}\ \ \texttt{>} & go up / down stairs\\
\texttt{o}\ \ \texttt{c} & open / close a door\\
\texttt{\^{}D} & kick\\
\texttt{s} & search for hidden doors and traps\\
\texttt{:} & look at what's here\\
\texttt{;} & look at a distant square\\
\texttt{F}+dir & force an attack into that square\\
\end{tabular}

\vspace{0.9em}
\textbf{Using items}\par\vspace{0.2em}
\begin{tabular}{@{}ll@{}}
\texttt{,} & pick up\\
\texttt{d}\ \ \texttt{D} & drop one / several\\
\texttt{i} & list your inventory\\
\texttt{e} & eat\\
\texttt{q} & quaff (drink) a potion\\
\texttt{r} & read a scroll or spellbook\\
\texttt{z} & zap a wand\\
\texttt{Z} & cast a spell\\
\texttt{a} & apply a tool (horn, lamp, pick-axe)\\
\texttt{C} & name an item or monster\\
\texttt{\textbackslash} & list what you've discovered\\
\texttt{@} & toggle autopickup\\
\end{tabular}
\end{minipage}
\hfill
\begin{minipage}[t]{0.48\textwidth}
\textbf{Weapons and armor}\par\vspace{0.2em}
\begin{tabular}{@{}ll@{}}
\texttt{w} & wield a weapon\\
\texttt{x} & swap main and off-hand\\
\texttt{X} & two-weapon combat on/off\\
\texttt{t}\ \ \texttt{f} & throw / fire from quiver\\
\texttt{Q} & set the quiver\\
\texttt{W}\ \ \texttt{T} & wear / take off armor\\
\texttt{P}\ \ \texttt{R} & put on / remove ring or amulet\\
\texttt{E} & engrave (write Elbereth)\\
\texttt{a} & reach-attack with a polearm or lance\\
\end{tabular}

\vspace{0.9em}
\textbf{Extended and special}\par\vspace{0.2em}
\begin{tabular}{@{}ll@{}}
\texttt{\#pray} & pray to your god\\
\texttt{\#dip} & dip an item (fountains, holy water)\\
\texttt{\#offer} & sacrifice on an altar\\
\texttt{\#loot} & loot a container\\
\texttt{\#force} & force a lock\\
\texttt{\#enhance} & improve weapon and spell skills\\
\texttt{\#invoke} & use an artifact's power\\
\texttt{\#chat} & talk (shop price quotes)\\
\texttt{\#ride} & mount or dismount a steed\\
\texttt{\#annotate} & note something about this level\\
\texttt{p} & pay your shop bill\\
\texttt{\^{}O} & dungeon overview\\
\texttt{\^{}X} & your attributes and intrinsics\\
\end{tabular}
\end{minipage}
\par\vspace{1.4em}
{\centering Commands written with \texttt{\#} are typed out in full. Most actions then ask for a direction or an item. \texttt{Esc} cancels; \texttt{Ctrl+A} repeats your last command.\par}
\normalsize
\clearpage
% --- Dungeon Emergency Checklist (verso, faces the command reference) ---
\thispagestyle{empty}
{\centering{\large\bfseries Dungeon Emergency Checklist}\par\vspace{0.3em}
{\itshape The Mazes of Menace are a dangerous place.\\
Consult this handy checklist when you find yourself in peril.}\par}
\vspace{1.5em}
\footnotesize
\begin{center}
\begin{minipage}{0.82\textwidth}
\begin{itemize}[label=\(\square\),leftmargin=1.6em,itemsep=0.35em,topsep=0pt]
\item \textbf{Do you have a healing potion?} Drink one. A potion of extra healing, full healing, or even gain level will lift your HP in a pinch. (p.~\pageref*{potion-healing})
\item \textbf{Do you have a wand of death?} Zap (\texttt{z}) it at your enemy to vanquish them instantly. Alternately, a wand of sleep or a thrown potion of paralysis or sleeping can freeze a foe long enough for you to flee. (p.~\pageref*{wand-death})
\item \textbf{Have you written Elbereth on the ground?} Engrave it in the dust (\texttt{E}, then type \texttt{Elbereth}) and stand on it; most monsters won't dare attack. (p.~\pageref*{elbereth})
\item \textbf{Can you teleport?} Read a scroll of teleportation, or zap a wand of teleportation at yourself, and vanish from the fight. (p.~\pageref*{wand-teleportation})
\item \textbf{Can you dig down?} Zap a wand of digging straight down and drop through the floor that instant. (A pick-axe digs down too, but it takes many turns, so it won't save you mid-fight.) (p.~\pageref*{wand-digging})
\item \textbf{Are the stairs close?} Flee to them (\texttt{<} or \texttt{>}) and leave the level behind. Most monsters won't follow you up the stairs, though a few will. (p.~\pageref*{when-the-fight-goes-badly})
\item \textbf{Is a crowd closing in?} A ring of conflict turns the mob against itself; a scroll of scare monster sends them fleeing. (p.~\pageref*{scroll-scare-monster})
\item \textbf{Are you blinded, confused, or stunned?} Apply (\texttt{a}) a unicorn horn to clear it. (p.~\pageref*{unicorn-horn-interactions})
\item \textbf{Are you turning to stone?} Eat a lizard corpse or any acidic corpse (an acid blob will do), or pray, before the timer runs out. (p.~\pageref*{petrification-stoning})
\end{itemize}

\vspace{0.9em}
\textbf{Who else to turn to}
\begin{itemize}[label=\(\square\),leftmargin=1.6em,itemsep=0.3em,topsep=0.2em]
\item \textbf{Your faithful pet} is the bodyguard you walked in with. Keep it fed and close. (p.~\pageref*{making-friends})
\item \textbf{A temple priest} will uncurse your gear and sell protection for a donation. (p.~\pageref*{donating-to-priests})
\item \textbf{The Oracle} answers your questions, for a price, if you can reach her. (p.~\pageref*{the-oracle})
\end{itemize}

\vspace{0.9em}
\textbf{Whatever you do, do NOT}
\begin{itemize}[label=\(\times\),leftmargin=1.6em,itemsep=0.3em,topsep=0.2em]
\item pray again right after praying. Your god is not a vending machine.
\item read an unidentified scroll hoping it's teleport. It may be aggravate monster.
\item swing from on top of your own Elbereth. It defiles the ward.
\item eat a fresh corpse to ``heal.'' Corpses don't restore HP, and some bite back.
\end{itemize}

\vspace{1em}
{\centering\textbf{If all else fails, \texttt{\#pray}.}\ \ Your god may pull you back from death itself, but only about once every thousand turns. Save it for the true emergency. (p.~\pageref*{prayer})\par}
\end{minipage}
\end{center}
\vspace{1.6em}
{\centering\itshape Tear out this page and store it in your Bag of Holding.\par}
\normalsize
\clearpage

% --- Title page ---
\begin{titlepage}
\thispagestyle{empty}
\vspace*{2in}
\begin{center}
{\huge\bfseries A Traveler's Companion to the\\[0.3em]Mazes of Menace}
\end{center}
\end{titlepage}

% --- Epigraph page ---
\thispagestyle{empty}
\vspace*{\fill}
\begin{center}
{\itshape Most of what follows is true.}
\end{center}
\vspace*{\fill}
\clearpage

% --- Table of Contents ---
\frontmatter
\thispagestyle{empty}
{
\hypersetup{linkcolor=black}
\setcounter{tocdepth}{0}
\small
\setlength{\parskip}{0.3em}
\renewcommand{\baselinestretch}{0.92}\selectfont
% Indent chapter entries in TOC to bring them closer to page numbers
\makeatletter
\renewcommand{\l@chapter}{\@dottedtocline{0}{1.5em}{0em}}
\makeatother
\tableofcontents
}

% --- Main content ---
\mainmatter

$body$

\end{document}
